\chapter{Product Thesis}
\label{ch:thesis}

\section{What Xyra is}

Xyra is a desktop integrated development environment whose primary user is a developer
working alongside autonomous agents. Every other editor treats the agent as a plugin
bolted onto a text editing surface designed for a human typing character by character.
Xyra inverts the relationship. The agent is a first class runtime with its own lifecycle,
its own verification loop, its own memory, and its own presence in the interface. The text
editor is one surface among several, and it is the surface a developer uses to inspect,
correct and approve work that mostly originates elsewhere.

The product exists because three assumptions behind current agentic editors are wrong.

\begin{enumerate}
  \item \req{Agents present unverified code.} A model writes a diff and the diff appears.
  Whether it compiles, whether the tests pass, whether it breaks a caller three modules
  away, all of that is discovered later by the human. Xyra refuses to display a change
  that has not been executed.

  \item \req{Context is treated as a bucket.} Editors stuff files into a prompt until the
  window fills, then truncate. This wastes the majority of the budget on tokens that carry
  no decision weight. Xyra treats context as a compression problem with a measurable
  objective and spends its budget on the smallest representation that preserves the ability
  to act correctly.

  \item \req{One session equals one task.} Work stops when the window fills, and the human
  restarts it. Xyra decomposes an objective into a dependency graph and runs each node in a
  fresh session, so total work is unbounded by any single context window.
\end{enumerate}

\section{The one sentence}

Xyra is the IDE where an autonomous team of agents ships a project while you sleep, and
every line they wrote was proven before you ever saw it.

\section{Design principles}

\begin{description}[style=nextline,leftmargin=0pt]
  \item[Proof before presentation.] No diff reaches a human surface until it has been
  executed against the project's own test or build command inside an isolated snapshot.
  This is enforced by the runtime, not by prompting.

  \item[Determinism where determinism is possible.] Import graphs, call graphs, symbol
  tables and dependency closures are computed, not inferred. Probabilistic retrieval is a
  fallback for questions that cannot be answered structurally, never the first resort.

  \item[Local by default.] Indexing, embedding, graph construction, compression, screenshot
  analysis, quality assurance runs and the entire mission state live on the developer's
  machine. Network calls are limited to the model provider the developer chose.

  \item[Every automated act is reversible and attributable.] One commit per unit of work,
  a decision journal that records why, and a time travel surface that can branch from any
  past decision.

  \item[The interface exposes the machine's state, not a chat transcript.] What the agent
  is reading, what it decided, what it verified, what it is about to do, and what it cost.
  A transcript is the least informative possible view of an autonomous system.

  \item[Nothing blocks on a human unless a human must decide.] The supervisor asks nothing
  it can resolve. When a decision genuinely belongs to the developer, it arrives as a
  structured choice with costs attached, never as a terminal question.
\end{description}

\section{Positioning}

\begin{table}[h]
\centering
\small
\begin{tabularx}{\textwidth}{l X X}
\toprule
\textbf{Capability} & \textbf{Mainstream agentic editors} & \textbf{Xyra} \\
\midrule
Verification of generated code & Human runs it after the fact & Runtime executes it in an isolated snapshot before display \\
Context assembly & Similarity search over chunks, truncation at the limit & Structural retrieval plus a compression pipeline with a token budget solver \\
Long horizon work & Bounded by one context window & Unbounded, ticket graph across fresh sessions \\
Visual correctness & Not modelled & Headless render, screenshot, vision verdict in the fix loop \\
Cross repository change & Manual, one repository at a time & Fleet impact analysis and sequential execution across repositories \\
Model relationship & Single vendor, metered & Multi vendor, adversarial review, flat rate capable \\
Failure handling & Session ends, human restarts & Retry, vendor swap, ticket split, quarantine, continue \\
State durability & Ephemeral chat & Durable mission state, decision journal, resumable after crash \\
\bottomrule
\end{tabularx}
\caption{Positioning against the current generation of agentic editors.}
\end{table}

\section{Why a standalone product rather than a fork}

The existing Xyra ships as a branded build of a GPL licensed editor. That path validated the
product thesis quickly and it remains a legitimate distribution channel, but it imposes two
permanent constraints. Source disclosure obligations follow every binary, which removes the
option of a proprietary commercial edition. Upstream churn dictates the rebase schedule,
which means engineering capacity is consumed by tracking another organisation's decisions.

The standalone product removes both constraints by building on permissively licensed
foundations only. Chapter~\ref{ch:foundation} enumerates the verified license position of
every dependency in the proposed stack. The strategic consequence is that the same codebase
can be released as a free edition and licensed as a commercial edition without any
restructuring.

\section{Product surfaces}

Xyra is one application with six primary surfaces. Each is specified control by control in
Chapter~\ref{ch:interface}.

\begin{table}[h]
\centering
\small
\begin{tabularx}{\textwidth}{l X}
\toprule
\textbf{Surface} & \textbf{Purpose} \\
\midrule
Workbench & Editing, navigation, diffing, terminal, source control. The classic IDE surface. \\
Mission Control & The autonomous supervisor: objective, ticket graph, progress, budget, controls. \\
Agent Theatre & Live orchestration of the agent team, lane by lane, with the artefacts each produced. \\
Context Inspector & What the agent is reading, weighted, with the ability to exclude and pin. \\
Verification Deck & Sandbox runs, visual checks, quality assurance sweeps, security and cost flags. \\
Atlas & Structural view of the codebase and the fleet: topology, impact, ownership, history. \\
\bottomrule
\end{tabularx}
\caption{The six product surfaces.}
\end{table}

\begin{figure}[h]
\centering
\begin{tikzpicture}[
  node distance=6mm,
  surf/.style={draw=xyraline,fill=xyrapanel,rounded corners=2pt,minimum height=9mm,minimum width=32mm,font=\small},
  core/.style={draw=xyralime,line width=1pt,fill=white,rounded corners=2pt,minimum height=11mm,minimum width=44mm,font=\small\bfseries}
]
\node[core] (rt) {Xyra Runtime};
\node[surf,above left=8mm and 14mm of rt] (wb) {Workbench};
\node[surf,above=8mm of rt] (mc) {Mission Control};
\node[surf,above right=8mm and 14mm of rt] (at) {Agent Theatre};
\node[surf,below left=8mm and 14mm of rt] (ci) {Context Inspector};
\node[surf,below=8mm of rt] (vd) {Verification Deck};
\node[surf,below right=8mm and 14mm of rt] (al) {Atlas};
\foreach \n in {wb,mc,at,ci,vd,al} \draw[xyragrey,-{Stealth[length=2mm]}] (rt) -- (\n);
\end{tikzpicture}
\caption{The runtime owns all state. Surfaces are projections of that state, never owners of it.}
\end{figure}

\section{Non goals}

Stating what Xyra does not attempt is as load bearing as stating what it does.

\begin{itemize}
  \item Xyra is not a general purpose text editor competing on startup time for editing a
  single configuration file. It is a project scale environment.
  \item Xyra does not host models. It orchestrates model providers the developer already
  pays for, including local ones.
  \item Xyra does not attempt to replace the shell or the browser. It drives them and reads
  their output.
  \item Xyra does not ship an interactive debugger in the first product generation. Runtime
  evidence reaches the agent through test output, traces and logs rather than through a
  breakpoint session, which is why no debug adapter appears in the issue register and no
  breakpoint control appears in the interface specification.
  \item Xyra does not implement its own version control. It drives git.
  \item Xyra does not implement real time collaborative editing. Multiple developers work
  through version control, and the collaboration surface in the title bar is reserved for a
  later generation.
  \item Xyra does not target mobile or web delivery in the first product generation.
\end{itemize}

The commercial model, pricing and edition split are deliberately outside the scope of this
document. Chapter~\ref{ch:foundation} establishes the licensing position that keeps every
commercial option open, which is the only part of that decision that constrains engineering.
